\section{Introduction}

Bipolar Disorder (BD) affects approximately 1 to 2\% of the global population \cite{merikangas2007prevalence}, yet epidemiological studies consistently find that patients wait nearly a decade before receiving a correct diagnosis \cite{hirschfeld2003perceptions}. The core reason is structural. Depressed patients seek clinical help, while hypomanic and manic episodes frequently go unreported or are experienced as periods of elevated productivity. A clinician encountering a patient in the middle of a depressive episode therefore has no direct behavioral evidence to separate BD from ordinary Major Depressive Disorder (MDD). This diagnostic confusion carries serious clinical consequences: antidepressant monotherapy prescribed to an unrecognized bipolar patient can trigger manic switching, rapid cycling, and mixed states, all of which worsen long-term prognosis \cite{ghaemi2004antidepressants}.

Current diagnostic practice relies on structured clinical interviews, self-report questionnaires such as the PHQ-9, MADRS, and HCL-32, and careful longitudinal observation. All of these approaches share a common vulnerability: they depend on the patient's retrospective account of their own mood states. Recall bias is well-documented in this population, and the cognitive impairments associated with depressive episodes themselves further compromise accurate self-reporting. There is therefore a compelling clinical motivation for objective, passive monitoring tools that operate independently of patient recall.

Wrist-worn actigraphy offers one such tool. Modern consumer-grade wearables continuously record movement-derived activity counts at minute-level resolution, generating multi-day behavioral time series without any active participation from the wearer. The psychiatric literature has long established that both bipolar and unipolar depression are accompanied by characteristic disruptions to rest-activity cycles and circadian rhythms \cite{wehr1987sleep,boudebesse2014sleep}. If these disruptions carry distinct temporal signatures across the two conditions, an automated classifier operating on raw actigraphy data could in principle assist clinicians in making the differential diagnosis earlier and more reliably.

Prior computational work on this problem has largely operated within the classical machine learning paradigm, applying support vector machines, random forests, and related methods to manually extracted statistical summaries of actigraphy signals \cite{garcia2018depresjon,zanella2019feature}. These features capture average activity levels reasonably well but miss subtler temporal structure: fragmented sleep bouts, irregular rest-activity transitions across days, and multi-day rhythmic drift. We argue that a deep sequence model that processes the raw time series directly, without prior feature engineering, is better positioned to discover the kinds of multi-scale temporal patterns that distinguish the two conditions \cite{lecun2015deep}.

In this paper, we ask a concrete empirical question: can a deep sequence model learn to distinguish bipolar from unipolar depression directly from raw wrist actigraphy, without manual feature engineering. We propose a hybrid 1D-CNN-LSTM architecture that pairs convolutional layers for localized kinetic pattern extraction with a recurrent LSTM module for multi-day information integration. We evaluate this architecture on the publicly available Depresjon dataset \cite{garcia2018depresjon} through a systematic four-approach investigation that progressively probes the signal from multiple angles: deep learning with synthetic oversampling, balanced training via downsampling, alternative neural architectures, and a diagnostic phase combining classical machine learning with direct statistical hypothesis testing.

The remainder of this paper is structured as follows. Section~\ref{sec:related} surveys related work across three domains: motor activity as a psychiatric biomarker, deep learning for wearable time-series, and the psychiatric literature on circadian disruption in mood disorders. Section~\ref{sec:methods} describes the dataset, preprocessing pipeline, and base model architecture. Section~\ref{sec:experiments} presents all four experimental approaches with full results and analysis. Section~\ref{sec:conclusions} draws conclusions, addresses open questions, and discusses the ethical implications of deploying automated diagnostic tools in mental health contexts.
